\section{Related Work}

\paragraph{CVSS vector prediction.}
\citet{elbaz2020fighting} predicted each CVSS metric from a bag-of-words description with linear models and derived the score from the predicted vector, with an emphasis on explainable decisions at disclosure time, and CVSS-BERT~\citep{shahid2021cvssbert} trained one BERT classifier per CVSS v3.1 metric and computed the score with the equations of the specification.
\citet{marchiori2025canllms} compared general-purpose LLMs with embedding-based classifiers for computing CVSS vectors and found the embedding methods stronger on the more subjective impact metrics of confidentiality, integrity, and availability.
The present study follows the predict-the-vector design of these works and adds a fine-tuned small language model, a strict temporal split, and a direct comparison with score regression.

\paragraph{Severity and CWE classification.}
VLAI~\citep{bonhomme2025vlai} fine-tunes RoBERTa on more than 600,000 vulnerabilities to predict the severity tier, but its training set spans CVEs from many years including 2024, which leaves it without a fair comparison on the 2024 test set used here (Section~\ref{sec:data}).
For CWE assignment, \citet{aota2020automation} used bag-of-words features with random forests, and \citet{mosievskiy2026roberta} reports that a 125M-parameter RoBERTa is competitive with much larger LLMs.

\paragraph{LLMs for vulnerability triage.}
\citet{talibzade2025llms} found that TF-IDF classifiers outperform general-purpose LLMs with generic prompts on CWE and severity assignment and attribute this to the schematic wording of CVE descriptions, while CVE-LLM~\citep{ghosh2025cvellm} fine-tunes an LLM with a domain ontology to evaluate medical-device vulnerabilities and \citet{torkamani2025casey} use LLMs to assign CWEs and severity in a triage tool.
For the related task of predicting exploitation, \citet{roytman2025benchmarking} reports that frontier LLMs fall well short of the statistical EPSS model and cost far more to run at scale, and industry systems such as VulnWatch~\citep{databricks2025vulnwatch} prioritize vulnerabilities with LLMs on hosted infrastructure.
The setting studied here differs from these systems in running entirely on local hardware and in evaluating every output on vulnerabilities disclosed after the training data.
